\section{Results}

Exploratory PPO results are available for the aligned 2D open-bay condition
(campaign \texttt{20261125-2331\_ppo\_ppo-ui-campaign}). Seed 0 completed
200,704 environment steps with 10 successes in 416 sampled training episodes.
Deterministic evaluation repeated the aligned start 20 times and exited the
workspace each time. The final replay crossed the target longitudinal
position with $0.088\,\mathrm{m}$ lateral error and $5.41^\circ$ heading
error, and missed the nominal heading and full-trailer containment criteria.
Peak articulation declined from $45.36^\circ$ at the 20,000-step checkpoint to
$8.38^\circ$ at the final checkpoint. This is exploratory evidence of more
stable reversing from the aligned start. A formal method comparison remains
pending.

On the single-direction quarter-arc pilot
(\path{20261126-1304_ppo_arc-one-direction-campaign}), guided PPO
completed 400,384 steps with zero docks in 625 training episodes.
Most training failures were jackknifes (542 episodes). The final policy
exited the workspace on all 20 validation starts, with mean arc coverage
of 14.8\% before exceeding a $0.5\,\mathrm{m}$ lateral-error threshold
or terminating. The 140,000-step checkpoint covered 19.1\% under the
same metric. LQR docked on all 20 matched starts within the 2,000-step
cap. This exploratory run does not isolate the effect
of removing mirrored starts because the earlier mirrored pilots used
a shorter episode limit and smaller training budget.

The shorter one-sided arc campaign
\path{20261126-1316_ppo_ppo-ui-campaign} recorded zero docks in 421
training episodes. At the selected 120,000-step checkpoint, all 20
validation replays completed the arc within $0.5\,\mathrm{m}$ of the
reference, with mean absolute arc error of $0.027\,\mathrm{m}$. All
20 subsequently exited the workspace after missing the docking pose.
The final checkpoint averaged 80.8\% corridor coverage and also failed
to dock in all 20 validation replays.

The exploratory gated-progress fine-tune
\path{20261126-1327_ppo_ppo-ui-campaign} used the selected short-arc
checkpoint as initialization and completed 200,704 additional steps
at a learning rate of $3\times10^{-4}$. It recorded zero docks in 152
training episodes. All 20 final validation replays exited the workspace.
Mean straight-approach lateral error increased from $0.95\,\mathrm{m}$
at initialization to $12.85\,\mathrm{m}$ in the final policy. The
screened 140,000-step fine-tuned checkpoint reached $1.70\,\mathrm{m}$.
An ungated run with matched fine-tuning settings is required to isolate
the effect of the gate.

The recovery-reward arm
\path{20261126-1348_ppo_arc30-recovery-campaign} ran for 60,416 steps and
recorded zero docks in 52 sampled training episodes. The selected
20,000-step fine-tuned checkpoint had $0.93\,\mathrm{m}$ mean absolute
straight-approach error, compared with $0.95\,\mathrm{m}$ at initialization
and $5.78\,\mathrm{m}$ in the final policy. Each checkpoint exited the
workspace on all 20 validation starts. None returned to within
$0.3\,\mathrm{m}$ of the reference after exceeding $0.5\,\mathrm{m}$ on
the straight approach. The matched gated-only control
\path{20261126-1358_ppo_arc30-recovery-control-campaign} also docked
zero times in 20 validation trials. Its mean straight-approach error
was $2.26\,\mathrm{m}$ at 20,000 steps and $3.72\,\mathrm{m}$ for the
final policy, versus $0.93$ and $5.78\,\mathrm{m}$ with the recovery
term. Neither condition recovered into the specified lateral corridor
on the straight leg at the checked checkpoints.

A weights-only warm start from the final aligned policy was trained on
small lateral and heading offsets against a straight centerline in the
same open bay. The validation result for
\path{20261125-2346_ppo-curriculum_ppo-open-dock-recovery-v1} was zero
docks and 20 workspace exits after 100,352 additional training steps.
The training episodes comprised 126 workspace exits and 48 jackknifes.
A paired diagnostic replay found that the mean absolute terminal lateral
error grew from $0.31\,\mathrm{m}$ for the initialization policy to
$4.56\,\mathrm{m}$ after recovery training. This outcome applies to the
single training seed and the validation scenario specified in the manifest.
A longer run under the same physical recovery condition
(\path{20261125-2359_ppo_ppo-ui-campaign}) finished with two docks in 347
sampled training episodes and jackknifed in all 20 deterministic validation
episodes after 200,704 training steps.

In the fixed gentle-curve campaign
\path{20261126-1213_ppo_ppo-ui-campaign}, PPO docked in 36 of 255 sampled
training episodes over 200,704 environment steps. The 100,000-step
checkpoint docked in 687 deterministic steps with $0.082\,\mathrm{m}$
mean absolute lateral path error on approach. The final checkpoint timed
out in all 20 repetitions of that fixed start. In the exploratory paired
comparison \path{20261126-1223_comparison_gentle-curve-nearby-diagnostic},
the selected PPO checkpoint docked in 15 of 20 nearby starts, while LQR
docked in 20 of 20. The PPO checkpoint was selected post hoc using replays
on the training curve. This provides local robustness evidence for one
training seed. The held-out comparison and path-guidance ablation remain
pending.

Engineering validation covers the shared core and simulator plants. The
recorded combined repository suite contains 454 passing tests
and four platform-dependent skips, including the 2D kinematic and MuJoCo
checks, controller vertical slices, and PPO training infrastructure, with
branch-aware coverage of $91.04\%$. Lint and strict
static type checks pass. The executable conformance report verifies
the neutral rigid-hitch constraint, tractor and trailer footprint areas,
integer 2D/MuJoCo timing ratios, nominal speed bounds, nominal trailer fit in
the configured dock, and semantic-site coverage. Kinematic checks exercise
forward and reverse propagation, constant steering, articulation response,
acceleration and steering limits, rigid reconstruction, deterministic replay,
and timestep convergence. Shared command sequences also compare the final 2D
and MuJoCo planar states within explicit test tolerances.

Gymnasium validation covers API conformance, seeded resets, fixed-reverse and
teleoperation actions, every terminal outcome, all three trailer lengths, a
1,000-step randomized stress rollout, and rendering. A constant straight
reverse from the default start reaches the configured success region after 486
control steps without collision. This is an environment acceptance test, not a
controller result.

The separately configured full-size service-center acceptance case verifies a
$14.0\,\mathrm{m}$ trailer, a $200\times100\,\mathrm{m}$ maneuvering workspace,
20 loading doors, absence of active-zone side-wall collision polygons, and
interactive motion up to $2\,\mathrm{m/s}$ forward and $1\,\mathrm{m/s}$
reverse. Boundary checks verify that articulation immediately below $60^\circ$
is nonterminal and articulation at $60^\circ$ is classified as jackknife. A
nonzero, unequal tractor/trailer-heading case also verifies that the rendered
tractor hitch, trailer hitch, and hitch marker remain coincident to numerical
precision after the dock-to-yard transform. Its constant straight reverse
reaches docking success after 325 control steps without collision. This remains
a visual and engineering demonstration, not an algorithm experiment.

The initial controller vertical-slice scenario uses the compact vehicle, an
explicit straight reference independent of the reset state, a $0.35\,\mathrm{m}$
lateral offset, a $3^\circ$ aligned tractor/trailer heading offset, fixed
$-0.20\,\mathrm{m/s}$ reverse speed, and the common steering-rate action. Both
the trajectory-error LQR and basic NMPC reach the configured docking success
event without collision or jackknife in deterministic acceptance runs. NMPC
reports no solver failure in this case. These checks establish end-to-end
controller integration only; they are one tuned scenario, are not a fair
performance comparison, and provide no evidence yet for curved paths,
randomized starts, or obstacle-constrained prediction.

These engineering checks establish internal consistency and low-speed model
behavior; they do not establish real-vehicle physical fidelity or controller
performance.
Numerical controller results, plots, confidence intervals, and failure
categories will be added only from saved experiment IDs.
